\documentclass[12pt]{article}
\usepackage[utf8]{inputenc}
\usepackage[catalan]{babel}
\usepackage{amsmath, amssymb, amsfonts}
\usepackage{geometry}
\usepackage{enumitem}
\geometry{a4paper, left=2cm, right=2cm, top=2.5cm, bottom=2.5cm}

\begin{document}

\begin{center}
    {\Large \textbf{FITXA DE DERIVADES}} \\
    \vspace{0.2cm}
    \textsc{Matemàtiques 2n Batxillerat}
\end{center}

\vspace{0.5cm}

\section*{Exercicis de càlcul de derivades}

\begin{enumerate}[label=\arabic*.]
    \item Calcula la derivada de les següents funcions utilitzant la definició (límit):
    \begin{enumerate}[label=\alph*)]
        \item $f(x)=3x+2$
        \item $f(x)=x^2-1$
        \item $f(x)=\dfrac{1}{x}$ \quad ($x\neq 0$)
    \end{enumerate}

    \item Calcula la derivada de les funcions següents:
    \begin{enumerate}[label=\alph*)]
        \item $f(x)=5x^7-3x^4+2x^2-9$
        \item $f(x)=\dfrac{x^2-1}{x^2+1}$
        \item $f(x)=e^x\cdot \ln x$
        \item $f(x)=\sin(2x+3)$
        \item $f(x)=\ln(x^2+1)$
        \item $f(x)=\cos^2 x$
        \item $f(x)=\sqrt{x^2+1}$
        \item $f(x)=\arctan(3x)$
    \end{enumerate}

    \item Troba l'equació de la recta tangent a les següents funcions en el punt indicat:
    \begin{enumerate}[label=\alph*)]
        \item $f(x)=x^2-4x+3$ en $x=2$
        \item $f(x)=\sqrt{x}$ en $x=4$
        \item $f(x)=e^x$ en $x=0$
        \item $f(x)=\ln x$ en $x=e$
    \end{enumerate}

    \item Calcula $f''(x)$ i $f'''(x)$ per a:
    \begin{enumerate}[label=\alph*)]
        \item $f(x)=x^5-2x^3+7x$
        \item $f(x)=\ln x$
        \item $f(x)=e^{2x}$
        \item $f(x)=\sin(2x)$
    \end{enumerate}

    \item Estudia la derivabilitat de les funcions següents en els punts indicats:
    \begin{enumerate}[label=\alph*)]
        \item $f(x)=|x|$ en $x=0$ (calcula els límits laterals del quocient incremental).
        \item $f(x)=
        \begin{cases}
            x^2 & \text{si } x\le 1 \\
            2x-1 & \text{si } x>1
        \end{cases}$
        en $x=1$.
    \end{enumerate}
\end{enumerate}

\section*{Aplicacions de les derivades}

\begin{enumerate}[label=\arabic*.]
    \setcounter{enumi}{5}
    \item Estudia la monotonia i els extrems relatius de les funcions:
    \begin{enumerate}[label=\alph*)]
        \item $f(x)=x^2-4x+3$
        \item $f(x)=x^3-3x^2+2$
        \item $f(x)=\dfrac{x}{x^2+1}$
    \end{enumerate}

    \item Determina els intervals de concavitat i els punts d'inflexió de:
    \begin{enumerate}[label=\alph*)]
        \item $f(x)=x^4-2x^2$
        \item $f(x)=\ln(x^2+1)$
        \item $f(x)=\sin x$ en $[0,2\pi]$
    \end{enumerate}

    \item Representa gràficament les funcions seguint el procediment complet (domini, simetries, talls, asímptotes, monotonia, concavitat):
    \begin{enumerate}[label=\alph*)]
        \item $f(x)=\dfrac{x+1}{x-2}$
        \item $f(x)=x^3-3x$
    \end{enumerate}

    \item Resol els següents problemes d'optimització:
    \begin{enumerate}[label=\alph*)]
        \item Es vol construir un dipòsit cilíndric obert per la part superior amb capacitat per a $1000$ litres. Quines dimensions (radi i alçada) minimitzen la superfície de material utilitzat?
        \item Troba els punts de la paràbola $y=x^2$ que estan a distància mínima del punt $(0,2)$.
        \item Una empresa fabrica $x$ unitats d'un producte. El cost de producció és $C(x)=3x^2-5x+12$ i el preu de venda és $P(x)=80-2x$. Troba el nombre d'unitats que maximitza el benefici.
    \end{enumerate}

    \item Calcula els següents límits utilitzant la regla de l'Hôpital:
    \begin{enumerate}[label=\alph*)]
        \item $\displaystyle \lim_{x\to 0} \dfrac{e^x-1}{x}$
        \item $\displaystyle \lim_{x\to 1} \dfrac{\ln x}{x-1}$
        \item $\displaystyle \lim_{x\to +\infty} \dfrac{\ln x}{x}$
        \item $\displaystyle \lim_{x\to 0^+} x^x$ (Indicació: pren logaritmes)
    \end{enumerate}
\end{enumerate}

\newpage

\section*{Solucionari}

\subsection*{Exercici 1}
\begin{enumerate}[label=\alph*)]
    \item $f'(x)=3$
    \item $f'(x)=2x$
    \item $f'(x)=-\dfrac{1}{x^2}$
\end{enumerate}

\subsection*{Exercici 2}
\begin{enumerate}[label=\alph*)]
    \item $f'(x)=35x^6-12x^3+4x$
    \item $f'(x)=\dfrac{4x}{(x^2+1)^2}$
    \item $f'(x)=e^x\left(\ln x+\dfrac{1}{x}\right)$
    \item $f'(x)=2\cos(2x+3)$
    \item $f'(x)=\dfrac{2x}{x^2+1}$
    \item $f'(x)=-2\cos x\sin x = -\sin(2x)$
    \item $f'(x)=\dfrac{x}{\sqrt{x^2+1}}$
    \item $f'(x)=\dfrac{3}{1+9x^2}$
\end{enumerate}

\subsection*{Exercici 3}
\begin{enumerate}[label=\alph*)]
    \item $y=-1$  (ja que $f(2)=-1$, $f'(2)=0$)
    \item $y=\frac{1}{4}x+1$ (ja que $f(4)=2$, $f'(4)=\frac{1}{4}$)
    \item $y=x+1$
    \item $y=\frac{1}{e}x$ (ja que $f(e)=1$, $f'(e)=1/e$)
\end{enumerate}

\subsection*{Exercici 4}
\begin{enumerate}[label=\alph*)]
    \item $f''(x)=20x^3-12x$, $f'''(x)=60x^2-12$
    \item $f''(x)=-\dfrac{1}{x^2}$, $f'''(x)=\dfrac{2}{x^3}$
    \item $f''(x)=4e^{2x}$, $f'''(x)=8e^{2x}$
    \item $f''(x)=-4\sin(2x)$, $f'''(x)=-8\cos(2x)$
\end{enumerate}

\subsection*{Exercici 5}
\begin{enumerate}[label=\alph*)]
    \item No derivable en $x=0$ (derivada per l'esquerra $-1$, per la dreta $+1$).
    \item Sí derivable en $x=1$ (ambdues derivades laterals valen $2$).
\end{enumerate}

\subsection*{Exercici 6}
\begin{enumerate}[label=\alph*)]
    \item $f'(x)=2x-4$. Decreix en $(-\infty,2)$, creix en $(2,+\infty)$. Mínim a $(2,-1)$.
    \item $f'(x)=3x(x-2)$. Creix en $(-\infty,0)$ i $(2,+\infty)$, decreix en $(0,2)$. Màxim a $(0,2)$, mínim a $(2,-2)$.
    \item $f'(x)=\dfrac{1-x^2}{(x^2+1)^2}$. Creix en $(-1,1)$, decreix en $(-\infty,-1)$ i $(1,+\infty)$. Mínim a $(-1,-1/2)$, màxim a $(1,1/2)$.
\end{enumerate}

\subsection*{Exercici 7}
\begin{enumerate}[label=\alph*)]
    \item $f''(x)=12x^2-4$. Punts d'inflexió: $x=\pm\frac{1}{\sqrt{3}}$. Còncava cap amunt en $(-\infty,-1/\sqrt{3})\cup(1/\sqrt{3},+\infty)$, cap avall en $(-1/\sqrt{3},1/\sqrt{3})$.
    \item $f''(x)=\dfrac{2(1-x^2)}{(x^2+1)^2}$. Inflexió a $x=\pm1$. Còncava amunt en $(-1,1)$, avall en els altres.
    \item $f''(x)=-\sin x$. Inflexió a $0,\pi,2\pi$. Còncava amunt en $(\pi,2\pi)$, avall en $(0,\pi)$.
\end{enumerate}

\subsection*{Exercici 8 (resum)}
\begin{enumerate}[label=\alph*)]
    \item Domini $\mathbb{R}\setminus\{2\}$, asímptota vertical $x=2$, obliqua $y=x+2$, màxim a $(0,-1/2)$, mínim a $(4,5/2)$.
    \item Domini $\mathbb{R}$, creixent/decreixent com a l'exercici 6(b), punts d'inflexió a $x=0$.
\end{enumerate}

\subsection*{Exercici 9}
\begin{enumerate}[label=\alph*)]
    \item $r=\sqrt[3]{\frac{1000}{\pi}}$ m, $h=2r$ (mínim). (En realitat, per a un cilindre obert superior, $h=r$ minimitza la superfície, però comprova-ho: $S=\pi r^2+2\pi r h$, $V=\pi r^2 h=1000$; la superfície mínima es dóna quan $h=r=\sqrt[3]{1000/\pi}$.)
    \item Punts $\left(\pm\frac{\sqrt{7}}{2},\frac{7}{4}\right)$.
    \item $x=20$ unitats (màxim de benefici $B(x)=P(x)x-C(x)$).
\end{enumerate}

\subsection*{Exercici 10}
\begin{enumerate}[label=\alph*)]
    \item $1$
    \item $1$
    \item $0$
    \item $1$ (prenent logaritmes: $\lim x\ln x=0$, per tant el límit és $e^0=1$).
\end{enumerate}



\newpage

\begin{center}
    {\Large \textbf{FITXA DE DERIVADES - NIVELL ALT}} \\
    \vspace{0.2cm}
    \textsc{Matemàtiques 2n Batxillerat (CCTT)}
\end{center}

\vspace{0.5cm}

\section*{Derivabilitat i continuïtat (amb paràmetres)}

\begin{enumerate}[label=\arabic*.]
    \item Troba els valors de $a$ i $b$ perquè la funció següent sigui derivable en $x=2$:
    \[
    f(x)=
    \begin{cases}
        ax^2 + bx + 1 & \text{si } x \le 2 \\
        3x - 1 & \text{si } x > 2
    \end{cases}
    \]

    \item Demostra que la funció $f(x)=|x^2-4|$ no és derivable en $x=-2$ ni en $x=2$, però sí ho és en la resta de punts.

    \item Sigui $f(x)=
    \begin{cases}
        x^2 \sin\left(\frac{1}{x}\right) & \text{si } x \neq 0 \\
        0 & \text{si } x = 0
    \end{cases}$.
    \begin{enumerate}[label=\alph*)]
        \item Demostra que $f$ és contínua en $x=0$.
        \item Demostra que $f$ és derivable en $x=0$ i calcula $f'(0)$.
        \item Demostra que $f'$ no és contínua en $x=0$. (Això demostra que la derivada d'una funció derivable no ha de ser contínua.)
    \end{enumerate}
\end{enumerate}

\section*{Càlcul de derivades (tècniques avançades)}

\begin{enumerate}[label=\arabic*.]
    \setcounter{enumi}{3}
    \item Calcula la derivada de les funcions següents utilitzant derivació logarítmica:
    \begin{enumerate}[label=\alph*)]
        \item $y = x^x$
        \item $y = \dfrac{(x+1)^3}{\sqrt{x-2}}$
        \item $y = (\sin x)^{\cos x}$
    \end{enumerate}

    \item Troba la derivada de les funcions implícites següents:
    \begin{enumerate}[label=\alph*)]
        \item $x^2 + y^2 = 25$
        \item $x^3 + y^3 = 6xy$ \quad (foli de Descartes)
        \item $\ln(x+y) = x$
    \end{enumerate}

    \item Calcula les derivades successives (fins a l'ordre $n$ que s'indica):
    \begin{enumerate}[label=\alph*)]
        \item $f(x)=\dfrac{1}{x}$. Troba $f^{(n)}(x)$.
        \item $f(x)=\ln(1+x)$. Troba $f^{(n)}(x)$.
        \item $f(x)=\sin(ax+b)$. Troba $f^{(n)}(x)$.
    \end{enumerate}
\end{enumerate}

\section*{Aplicacions teòriques i teoremes}

\begin{enumerate}[label=\arabic*.]
    \setcounter{enumi}{6}
    \item Aplica el teorema de Rolle per demostrar que l'equació $2x^3 - 3x^2 + 1 = 0$ té, com a mínim, dues arrels reals en l'interval $(-1,2)$.

    \item Demostra que l'equació $x^3 + 3x - 1 = 0$ té exactament una arrel real. (Indicació: primer demostra que existeix per Bolzano, després que és única per monotonia estricta.)

    \item Sigui $f$ contínua en $[a,b]$ i derivable en $(a,b)$ amb $f(a)=f(b)=0$. Demostra que per a qualsevol $k \in \mathbb{R}$, existeix $c \in (a,b)$ tal que:
    \[
    f'(c) = k f(c)
    \]
    (Indicació: aplica Rolle a $g(x)=e^{-kx}f(x)$.)

    \item Un punt es mou segons la llei $s(t) = t^3 - 6t^2 + 9t + 2$ (en metres). Troba:
    \begin{enumerate}[label=\alph*)]
        \item La velocitat i l'acceleració en funció del temps.
        \item En quins intervals la velocitat és positiva.
        \item L'acceleració en els instants en què la velocitat és nul·la.
    \end{enumerate}
\end{enumerate}

\section*{Optimització i geometria}

\begin{enumerate}[label=\arabic*.]
    \setcounter{enumi}{10}
    \item Es vol construir un dipòsit cilíndric tancat de volum $V$ (constant). Demostra que la superfície total és mínima quan l'alçada és igual al diàmetre.

    \item Una finestra té forma de rectangle coronat per un semicercle. El perímetre total és de $10$ m. Troba les dimensions perquè l'àrea de la finestra sigui màxima.

    \item Troba el punt de la corba $y^2 = 4x$ que està més a prop del punt $(4,0)$.

    \item Sobre la circumferència $x^2 + y^2 = 1$, determina els punts on la tangent és paral·lela a la recta $y = x + 2$.
\end{enumerate}

\section*{Regla de l'Hôpital (casos complexos)}

\begin{enumerate}[label=\arabic*.]
    \setcounter{enumi}{14}
    \item Calcula els límits següents (transforma prèviament la indeterminació):
    \begin{enumerate}[label=\alph*)]
        \item $\displaystyle \lim_{x\to 0^+} x^{\sin x}$
        \item $\displaystyle \lim_{x\to \infty} \left(1+\frac{1}{x^2}\right)^x$
        \item $\displaystyle \lim_{x\to 0} \left(\frac{1}{\sin x} - \frac{1}{x}\right)$
        \item $\displaystyle \lim_{x\to 1} \left(\frac{1}{\ln x} - \frac{1}{x-1}\right)$
    \end{enumerate}
\end{enumerate}

\newpage

\section*{Solucionari (amb indicacions)}

\begin{enumerate}[label=\arabic*.]
    \item Per continuïtat: $\lim_{x\to 2^-} f(x)=f(2)=4a+2b+1$. Límit dreta: $5$. Per tant $4a+2b=4$ (i).
    Per derivabilitat: derivada esquerra $=4a+b$, dreta $=3$. Per tant $4a+b=3$ (ii).
    Resolent: $a= \frac{1}{2}$, $b=1$.

    \item Com que $|x^2-4| = |(x-2)(x+2)|$, el valor absolut crea un "pic" a $x=\pm2$ (on el signe canvia). Les derivades laterals no coincideixen (valen $\pm 4$), per tant no és derivable. En la resta, la funció és un polinomi amb signe constant, per tant derivable.

    \item 
    \begin{enumerate}[label=\alph*)]
        \item $0 \le |f(x)| \le x^2 \to 0$, per tant contínua.
        \item $f'(0)=\lim_{h\to0} \frac{h^2\sin(1/h)}{h} = \lim h\sin(1/h)=0$.
        \item Per a $x\neq0$, $f'(x)=2x\sin(1/x)-\cos(1/x)$. Quan $x\to0$, $\cos(1/x)$ no té límit, per tant $f'$ no és contínua en $0$.
    \end{enumerate}

    \item 
    \begin{enumerate}[label=\alph*)]
        \item $\ln y = x\ln x \Rightarrow y' = x^x(\ln x+1)$
        \item $\ln y = 3\ln(x+1) - \frac{1}{2}\ln(x-2) \Rightarrow y' = \frac{(x+1)^3}{\sqrt{x-2}}\left(\frac{3}{x+1} - \frac{1}{2(x-2)}\right)$
        \item $\ln y = \cos x \ln(\sin x) \Rightarrow y' = (\sin x)^{\cos x}(-\sin x\ln(\sin x)+\cos x\cot x)$
    \end{enumerate}

    \item 
    \begin{enumerate}[label=\alph*)]
        \item $y' = -\frac{x}{y}$
        \item $3x^2+3y^2y'=6y+6xy' \Rightarrow y'=\frac{2y-x^2}{y^2-2x}$
        \item $1/(x+y)(1+y')=1 \Rightarrow y'=x+y-1$
    \end{enumerate}

    \item 
    \begin{enumerate}[label=\alph*)]
        \item $f^{(n)}(x)=(-1)^n \frac{n!}{x^{n+1}}$
        \item $f^{(n)}(x)=(-1)^{n-1}\frac{(n-1)!}{(1+x)^n}$
        \item Si $n$ parell: $(-1)^{n/2}a^n\sin(ax+b)$; si senar: $(-1)^{(n-1)/2}a^n\cos(ax+b)$.
    \end{enumerate}

    \item $f(-1)= -2-3+1=-4$, $f(0)=1$, $f(1)=0$, $f(2)=16-12+1=5$. Com que $f(-1)<0$, $f(0)>0$ $\Rightarrow$ arrel a $(-1,0)$; $f(0)>0$, $f(1)=0$ (ja n'hi ha una); $f(1)=0$, $f(2)>0$ $\Rightarrow$ arrel a $(1,2)$. Aplicant Rolle a cada interval es demostra l'existència de dues arrels, però cal veure que l'equació cúbica té com a màxim 3 arrels. N'hi ha prou amb veure que hi ha una arrel doble? En realitat, per Rolle, si tingués dues arrels diferents, la derivada s'anul·laria. $f'(x)=6x(x-1)$. Anul·lant a $0$ i $1$. Això demostra que hi ha com a mínim 2 arrels (una a $(-1,0)$ i una a $(1,2)$).
    \item Bolzano: $f(0)=-1<0$, $f(1)=3>0$ $\Rightarrow$ arrel. $f'(x)=3x^2+3>0$ sempre, per tant estrictament creixent, només una arrel.
    \item Definir $g(x)=e^{-kx}f(x)$. $g(a)=g(b)=0$. Rolle: existeix $c$ tal que $g'(c)=e^{-kc}(f'(c)-k f(c))=0$. Com que l'exponencial no s'anul·la, $f'(c)=k f(c)$.
    \item $v(t)=3t^2-12t+9$, $a(t)=6t-12$. $v>0$ per $t<1$ o $t>3$. $v=0$ a $t=1$ ($a=-6$) i $t=3$ ($a=6$).

    \item $S=2\pi r^2+2\pi r h$, $V=\pi r^2 h$ (constant). Aïllar $h=V/(\pi r^2)$, $S(r)=2\pi r^2+2V/r$. Derivar, igualar a zero: $4\pi r - 2V/r^2=0 \Rightarrow r^3=V/(2\pi)$. Llavors $h=V/(\pi r^2)=2r$, per tant $h=2r$ (diàmetre).

    \item Sigui $x$ l'amplada del rectangle, $y$ l'alçada. Radi semicercle $r=x/2$. Perímetre: $x+2y+\pi x/2 = 10$. Àrea: $A=xy+\frac{1}{2}\pi (x/2)^2 = xy+\pi x^2/8$. Aïllar $y$ i optimitzar. Resultat: $x = \frac{20}{4+\pi}$, $y = \frac{10}{4+\pi}$.

    \item Minimitzar $d^2=(x-4)^2+y^2$ subjecte a $y^2=4x$. Substituir: $d^2=(x-4)^2+4x = x^2-4x+16$. Derivar: $2x-4=0 \Rightarrow x=2$, $y=\pm 2\sqrt{2}$. Distància mínima $\sqrt{12}=2\sqrt{3}$.

    \item Pendent tangent $y' = -x/y$. Volem $y'=1 \Rightarrow -x/y=1 \Rightarrow y=-x$. Substituir a la circumferència: $2x^2=1 \Rightarrow x=\pm 1/\sqrt{2}$, $y=\mp 1/\sqrt{2}$. Punts: $(1/\sqrt{2}, -1/\sqrt{2})$ i $(-1/\sqrt{2}, 1/\sqrt{2})$.

    \item 
    \begin{enumerate}[label=\alph*)]
        \item $\lim x^{\sin x} = \exp(\lim \sin x \ln x) = \exp(\lim x \ln x)=e^0=1$.
        \item $\lim x \ln(1+1/x^2) = \lim x \cdot (1/x^2) = 0 \Rightarrow$ límit $=1$.
        \item $\lim \frac{x-\sin x}{x\sin x} = \lim \frac{1-\cos x}{\sin x + x\cos x} = \lim \frac{\sin x}{2\cos x - x\sin x} = 0$.
        \item $\lim \frac{x-1-\ln x}{(x-1)\ln x} = \lim \frac{1-1/x}{\ln x + (x-1)/x} = \lim \frac{1/x}{1/x + 1/x^2} = 1/2$.
    \end{enumerate}
\end{enumerate}

\end{document}